% ============================================================
% The Dynamical Substrate Law:
%   Substrate-Locality of DI-Bit Currents at Vertex-Aligned
%   Reading C in the 600-Cell
% Conscious Point Physics Flagship Paper Series (F-line)
%
% Version 0.1 -- 23 May 2026 (Session 142, Patch 0554 — Paper skeleton
%   creation. First .tex artifact for the F.1 flagship paper assembly
%   multi-Patch arc opened at Patch 0553. Preamble + section headers
%   + title + abstract draft only; no body content at this Patch.
%   Body sections to be added one-per-Patch across Patches 0555-0564
%   per F1_flagship_paper_assembly_scoping.md §4.2-§4.5. Three
%   publication-grade hardened theorem artifacts at
%   `flagship_papers/dynamical_substrate_law/hardened_theorems/`
%   (Patches 0550 + 0551 + 0552, 741 lines LaTeX combined) supply
%   §5 and §6 body content as direct integration inputs. ChatGPT
%   load-bearing principle codified at scoping doc anti-priority #8:
%   preserve Layer distinction + conditionality + open higher-order
%   questions; do not erase uncertainty structure during paper
%   polishing.)
% ============================================================

\documentclass[11pt,letterpaper]{article}
\usepackage[utf8]{inputenc}
\usepackage[margin=1in]{geometry}
\usepackage[T1]{fontenc}
\usepackage{amsmath,amssymb,amsfonts,amsthm}
\usepackage{graphicx}
\usepackage{tikz}
\usetikzlibrary{positioning,calc,arrows.meta,shapes.geometric}
\usepackage{xcolor}
\usepackage{hyperref}
\usepackage{booktabs}
\usepackage{array}
\usepackage{enumitem}
\usepackage{float}
\usepackage[font=small, labelfont=bf]{caption}
\usepackage{parskip}
\usepackage{mdframed}

\hypersetup{
    colorlinks=true,
    linkcolor=blue,
    citecolor=blue,
    urlcolor=blue
}

\newtheorem{theorem}{Theorem}[section]
\newtheorem{proposition}[theorem]{Proposition}
\newtheorem{lemma}[theorem]{Lemma}
\newtheorem{corollary}[theorem]{Corollary}
\newtheorem{conjecture}[theorem]{Conjecture}
\newtheorem{definition}[theorem]{Definition}
\newtheorem{remark}[theorem]{Remark}
\newtheorem{claim}[theorem]{Claim}
\newtheorem{openproblem}{Open Problem}

% Notation conventions inherited from Capotauro v2.0 + SF-4 +
% the F.1 hardened theorem artifacts (Patches 0550 + 0551 + 0552)
\newcommand{\phig}{\varphi}
\newcommand{\nhat}{\hat{n}}
\newcommand{\vhost}{v_{\text{host}}}
\newcommand{\jDInet}{\vec{j}_{DI}^{\,\text{net}}}
\newcommand{\omegaPCD}{\vec{\omega}_{PCD}}
\newcommand{\sigcycle}{\sigma_{\text{cycle}}}
\newcommand{\Mthermo}{|M^{\text{thermo}}|}
\newcommand{\DeltapLR}{\Delta p_{LR}}
\newcommand{\Ih}{I_h}
\newcommand{\PTI}{P_{TI}}
\newcommand{\Zbw}{\mathrm{ZBW}}
\newcommand{\bigO}[1]{\mathcal{O}(#1)}
\newcommand{\Findd}{\mathrm{DSL}}

\title{The Dynamical Substrate Law:\\
Substrate-Locality of DI-Bit Currents at\\
Vertex-Aligned Reading~C in the 600-Cell}

\author{Dr.~Thomas Lee Abshier, ND\\
\small Hyperphysics Institute\\
\small \texttt{drthomas@renaissanceministries.org}\\
\\
\small with AI co-authorship:\\
\small Claude Opus (Anthropic) --- primary mathematical co-author\\
\small ChatGPT (OpenAI), Copilot (Microsoft), Grok (xAI) --- review and calibration
}

\date{Version 0.1 (skeleton), 23 May 2026}

\begin{document}

\maketitle

% ============================================================
% Abstract --- draft v0.1, refined at later assembly Patches
% ============================================================
\begin{abstract}
We establish the substrate-locality property for DI-bit currents on the 600-cell substrate at vertex-aligned Reading~C: under Mechanism~A (the propagation-rate-asymmetry primitive $r(\hat{e}) = r_0(1 + \delta\,\hat{e}\cdot\nhat)$) and the framework-local current construction at first order in $\delta$, the substrate net DI-bit current at any host vertex $\vhost$ depends only on first-shell content. The result is proved at sketch-document Layer~3 in three steps: (i)~the host-to-first-shell unit-direction projection $\hat{u}_i \cdot \nhat = -1/(2\phig)$ is uniform across all twelve first-shell neighbours (Theorem~\ref{thm:host-first-shell-projection}); (ii)~the first-shell-to-first-shell edge-direction projection $\hat{e}_{ij} \cdot \nhat = 0$ vanishes for any pair of first-shell vertices (Theorem~\ref{thm:first-shell-perpendicularity}); (iii)~the perturbation-theory propagation rule confines the $\bigO{\delta^n}$ coefficient of the substrate current to graph-distance-$n$ edges in the 600-cell edge graph (Theorem~\ref{thm:perturbation-locality}, with shell-locality Corollary~\ref{cor:shell-locality}).
The three load-bearing identities are hardened at publication-grade rigor with explicit hypothesis tracking and five-class exclusion enumeration each.
We make the scope qualifier --- perturbative locality under Mechanism~A and the framework-local current construction at $\bigO{\delta^1}$ --- structurally inseparable from the theorem statements.
Three explicit open higher-order questions are flagged: extension to $\bigO{\delta^2}$ (Open~Problem~\ref{op:delta-squared}); Layer~4 axiomatic derivation of Mechanism~A from the eleven CPP primitive axioms~A1--A11 (Open~Problem~\ref{op:layer4-mechanism-a}); and publication-grade hardening of the first-shell inner-product primitive G1 (Open~Problem~\ref{op:g1-hardening}).
The paper integrates the F.1 sub-question of the OPEN-SD-CHIR-PRIMITIVE manifestation~(iv) trajectory and is the temporal-sector analog to the Capotauro~v2.0 spatial-sector substrate-locality theorem, with which it shares the structural constant $-1/(2\phig)$.
\end{abstract}

\tableofcontents

% ============================================================
% §1 Introduction and motivation
%   Body to be added at Patch 0555.
%   Primary source: F1_subquestion_pcd_orientation_link.md §1
%   (with framing for outer audience).
%   Approximate target: 3-4 pages.
% ============================================================
\section{Introduction and motivation}
\label{sec:introduction}

% [Body content to be added at Patch 0555 per scoping doc §3.2.
%  Sub-structure anticipated:
%   §1.1 Position in the CPP corpus and the chirality continuum
%   §1.2 The F.1 sub-question as a programme gate
%   §1.3 Sketch Layer 2 → Layer 3 trajectory recap
%   §1.4 Paper roadmap and Layer-distinction discipline]

% ============================================================
% §2 The F.1 sub-question (formal statement)
%   Body to be added at Patch 0556.
%   Primary source: F1_subquestion_pcd_orientation_link.md §2-§3
%   + F1_phase2_foundations_work.md §1-§3.
%   Approximate target: 2-3 pages.
% ============================================================
\section{The F.1 sub-question: substrate-mechanism derivation of the link
$\hat{n} \mapsto \omegaPCD$}
\label{sec:subquestion}

% [Body content to be added at Patch 0556. Sub-structure anticipated:
%   §2.1 Manifestation (iv) of OPEN-SD-CHIR-PRIMITIVE
%   §2.2 Three closure scenarios (A positive / B negative / C partial)
%   §2.3 Scenario A: Mechanism A propagation-rate asymmetry
%   §2.4 The minimal-local-first-order realization framework]

% ============================================================
% §3 Framework primitives
%   Body to be added at Patch 0557.
%   Primary source: F1_subquestion_pcd_orientation_link.md §4-§5
%   + axiom-registry A1/A6'/A11 + the 600-cell geometric structure
%   + Reading C from Capotauro v2.0.
%   Approximate target: 3-4 pages.
% ============================================================
\section{Framework primitives}
\label{sec:framework-primitives}

% [Body content to be added at Patch 0557. Sub-structure anticipated:
%   §3.1 CPP primitive axioms A1-A11 (recap)
%   §3.2 The 600-cell substrate at vertex-aligned Reading C
%   §3.3 First-shell geometric primitives G1 + G2
%   §3.4 The DI-bit current and Mechanism A framework]

% ============================================================
% §4 Mechanism A as framework axiom
%   Body to be added at Patch 0558.
%   Primary source: Patch 0544 §3 + Patch 0550 §2.2.
%   Approximate target: 2 pages.
% ============================================================
\section{Mechanism A as framework axiom}
\label{sec:mechanism-a}

% [Body content to be added at Patch 0558. Sub-structure anticipated:
%   §4.1 The propagation-rate asymmetry r(e) = r_0 (1 + delta * e.n)
%   §4.2 Framework local-current construction at first order in delta
%   §4.3 Mechanism A as framework axiom (Layer 3); Layer 4 derivation deferred]

% ============================================================
% §5 The first-shell geometric identities
%   Body to be added at Patch 0559.
%   INCORPORATES: hardened_theorems/first_shell_perpendicularity.tex
%   (Patch 0551, 207 lines) + host_to_first_shell_projection.tex
%   (Patch 0552, 265 lines). Integrated treatment of Theorems
%   4.1.1 + 4.2.1 as a single chapter with shared G1 dependency
%   (exclusion class E1).
%   Approximate target: 5-6 pages.
% ============================================================
\section{The first-shell geometric identities}
\label{sec:first-shell-identities}

% [Body content to be added at Patch 0559. Sub-structure anticipated:
%   §5.1 The G1 first-shell inner-product primitive (inherited Layer 3)
%   §5.2 Lemma 3.1.1: tangent-hyperplane chord length (1/phi)
%   §5.3 Theorem 5.1 (host-to-first-shell uniform projection;
%        from Patch 0552 .tex content)
%   §5.4 Theorem 5.2 (first-shell-to-first-shell perpendicularity;
%        from Patch 0551 .tex content)
%   §5.5 Exclusion class E1 (shared G1 dependency)
%   §5.6 Cross-reference: Capotauro v2.0 §3 spatial-sector parallel]
\begin{theorem}[Host-to-first-shell uniform projection; sketch-document Layer 3]
\label{thm:host-first-shell-projection}
[Statement to be added at Patch 0559; integrates Patch 0552 Theorem 4.2.1.]
\end{theorem}

\begin{theorem}[First-shell-to-first-shell perpendicularity; sketch-document Layer 3]
\label{thm:first-shell-perpendicularity}
[Statement to be added at Patch 0559; integrates Patch 0551 Theorem 4.1.1.]
\end{theorem}

% ============================================================
% §6 The perturbation-theory propagation rule
%   Body to be added at Patch 0560.
%   INCORPORATES: hardened_theorems/perturbation_locality_propagation.tex
%   (Patch 0550, 269 lines). Three lemmas (path-amplitude expansion +
%   perturbed-step counting + connected-subgraph confinement) +
%   shell-locality corollary + delta^1 specialisation.
%   Approximate target: 5-6 pages.
% ============================================================
\section{The perturbation-theory propagation rule}
\label{sec:perturbation-locality}

% [Body content to be added at Patch 0560. Sub-structure anticipated:
%   §6.1 Lemma: path-amplitude expansion
%   §6.2 Lemma: perturbed-step counting
%   §6.3 Lemma: connected-subgraph confinement
%   §6.4 Theorem 6.1: O(delta^n) coefficient confined to graph-distance-n edges
%   §6.5 Corollary 6.2: shell-locality at O(delta^1)
%   §6.6 Five-class exclusion enumeration]
\begin{theorem}[Perturbation-theory propagation rule; sketch-document Layer 3]
\label{thm:perturbation-locality}
[Statement to be added at Patch 0560; integrates Patch 0550 Theorem 3.3.3.]
\end{theorem}

\begin{corollary}[Shell-locality at $\bigO{\delta^1}$; sketch-document Layer 3]
\label{cor:shell-locality}
[Statement to be added at Patch 0560; integrates Patch 0550 Corollary 3.3.4.]
\end{corollary}

% ============================================================
% §7 The substrate-locality theorem (umbrella result)
%   Body to be added at Patch 0561.
%   Primary source: Patch 0544 §5 Theorem 5.1.1
%   (sketch-document Layer 3; uses Theorems 5.1, 5.2, 6.1 from §5+§6 as inputs).
%   Approximate target: 4-5 pages.
% ============================================================
\section{The substrate-locality theorem (umbrella result)}
\label{sec:substrate-locality-theorem}

% [Body content to be added at Patch 0561. Sub-structure anticipated:
%   §7.1 Theorem 7.1: substrate net DI-bit current at v_host depends
%        only on first-shell content at O(delta)
%   §7.2 Proof part (a): host-to-first-shell uniform contribution
%        (via Theorem 5.1)
%   §7.3 Proof part (b): first-shell-to-first-shell zero contribution
%        (via Theorem 5.2)
%   §7.4 Proof part (c): higher-shell zero contribution at O(delta)
%        (via Theorem 6.1 + Corollary 6.2)
%   §7.5 Connection to Phase 1 Finding DSL-1]
\begin{theorem}[Substrate-locality of DI-bit currents at vertex-aligned Reading C;
sketch-document Layer 3]
\label{thm:substrate-locality}
[Statement to be added at Patch 0561; integrates Patch 0544 §5 Theorem 5.1.1.]
\end{theorem}

% ============================================================
% §8 Layer 3 stack: status, scope, and exclusions
%   Body to be added at Patch 0562.
%   Primary source: Patch 0549 status framing + Patches 0550/0551/0552 §6
%   exclusion classes + per-claim Layer distinction labelling.
%   Approximate target: 3-4 pages.
% ============================================================
\section{Layer 3 stack: status, scope, and exclusions}
\label{sec:layer3-stack}

% [Body content to be added at Patch 0562. Sub-structure anticipated:
%   §8.1 The minimal-local-first-order realization framework
%   §8.2 Layer distinction labels: per-theorem rigor classification
%   §8.3 The five-class exclusion enumeration synthesis
%   §8.4 The Patch 0549 sub-question status framing recap
%   §8.5 Anti-erasure discipline (per ChatGPT Patch 0538 §10.7)]

% ============================================================
% §9 Open higher-order questions
%   Body to be added at Patch 0563.
%   Primary source: Patch 0549 §38.3 + the three .tex artifacts'
%   E-class lists.
%   Approximate target: 2-3 pages.
% ============================================================
\section{Open higher-order questions}
\label{sec:open-questions}

% [Body content to be added at Patch 0563. Sub-structure anticipated:
%   Five explicit open problems: OPEN-FP-F1-1 through OPEN-FP-F1-5.]
\begin{openproblem}[Extension to $\bigO{\delta^2}$; OPEN-FP-F1-1]
\label{op:delta-squared}
[Statement to be added at Patch 0563.]
\end{openproblem}

\begin{openproblem}[Layer 4 axiomatic derivation of Mechanism A; OPEN-FP-F1-2]
\label{op:layer4-mechanism-a}
[Statement to be added at Patch 0563.]
\end{openproblem}

\begin{openproblem}[Publication-grade hardening of identity G1; OPEN-FP-F1-3]
\label{op:g1-hardening}
[Statement to be added at Patch 0563.]
\end{openproblem}

\begin{openproblem}[Sector-5 schema instantiation; OPEN-FP-F1-4]
\label{op:sector5-schema}
[Statement to be added at Patch 0563.]
\end{openproblem}

\begin{openproblem}[Non-vertex-aligned Reading C variants; OPEN-FP-F1-5]
\label{op:non-vertex-aligned-c}
[Statement to be added at Patch 0563.]
\end{openproblem}

% ============================================================
% §10 Conclusion
%   Body to be added at Patch 0564.
%   Ties Layer 3 closure to CPP chirality-from-polytope-geometry
%   long-term programme target (Patch 0528 §14.17 + Patch 0539 §15.5).
%   Approximate target: 1-2 pages.
% ============================================================
\section{Conclusion}
\label{sec:conclusion}

% [Body content to be added at Patch 0564.]

% ============================================================
% Bibliography
%   Body to be added at Patch 0565.
%   Imports from bibliography/cpp_references.bib + adds F.1-specific
%   entries.
% ============================================================
\bibliographystyle{plain}
% \bibliography{../../bibliography/cpp_references}  % uncomment at Patch 0565
\begin{thebibliography}{99}

% [Bibliography entries to be added at Patch 0565.]

\end{thebibliography}

\end{document}
